\chapter*{Введение}
\addcontentsline{toc}{chapter}{Введение}
\label{ch:intro}

\textbf{Электронная почта} остаётся одним из ключевых инструментов деловой
коммуникации. Для корпоративных сетей необходим собственный почтовый сервер,
обеспечивающий конфиденциальность переписки, гибкое управление учётными записями
и независимость от внешних провайдеров.

Цель работы --- развернуть почтовый сервер на базе \textbf{iRedMail}
на виртуальной машине \texttt{mail} в локальной сети, созданной
в лабораторных работах~4 и~5.

Задачи:
\begin{enumerate}
  \item Подготовить виртуальную машину \texttt{mail}: назначить статический IP,
        настроить FQDN и DNS-записи.
  \item Установить iRedMail с бэкендом \textbf{OpenLDAP} и веб-сервером
        \textbf{Nginx}.
  \item Настроить почтовый домен \texttt{yazikov.iks531.local}.
  \item Создать почтовых пользователей через панель администратора \textbf{iRedAdmin}.
  \item Проверить отправку и получение писем через веб-клиент \textbf{Roundcubemail}.
\end{enumerate}

\textbf{Стенд:}
\begin{itemize}
  \item ВМ \texttt{gateway} --- Ubuntu~22.04~LTS (VirtualBox),
        два адаптера: \texttt{enp0s3}~(WAN/NAT, \texttt{10.0.2.15/24})
        и \texttt{enp0s8}~(LAN~Internal, \texttt{192.168.\cfgLabStudentN.1/24});
        роль --- шлюз, DNS-сервер (BIND9), DHCP-сервер.
  \item ВМ \texttt{mail} --- Ubuntu~22.04~LTS, один адаптер
        \texttt{enp0s3}~(Internal Network),
        IP \texttt{192.168.\cfgLabStudentN.5/24};
        роль --- почтовый сервер iRedMail.
  \item ВМ \texttt{Desktop} --- Ubuntu~22.04~Desktop, один адаптер
        (Internal Network), IP получает по DHCP;
        роль --- клиент для тестирования веб-почты.
\end{itemize}

\endinput
